\chapter{Sorrifier Examples and Case Studies}
\label{app:sorrifier-examples}

This subsection illustrates how the Sorrifier converts different Lean failures into
placeholder-compilable scripts and how the line-based syntactic reward is computed. In all examples,
blank lines and comments are excluded from the line count.

\paragraph{Case 1: Parser error.}
Consider a generated proof body containing a malformed tactic:
\begin{verbatim}
have h := 1
apply h (
\end{verbatim}
Lean reports a parser error because the opening parenthesis is not closed. The Sorrifier localizes
the failing tactic line and replaces it with \texttt{sorry}:
\begin{verbatim}
have h := 1
sorry
\end{verbatim}
The original body contains two clean lines. One line survives unchanged, while the malformed tactic is
replaced. Therefore, the syntactic reward is:
\[
r_{\text{env}} = \frac{1}{2} = 0.5.
\]

\paragraph{Case 2: Unresolved name and cascade cleanup.}
Consider a generated proof body that calls a nonexistent lemma:
\begin{verbatim}
apply nonexistent_lemma
rfl
\end{verbatim}
Lean reports an unknown identifier error for \texttt{nonexistent\_lemma}. After this failure, the following
line is no longer useful in the recovered proof context. The Sorrifier therefore replaces the failed
region with a single placeholder:
\begin{verbatim}
sorry
\end{verbatim}
Although the original body has two clean lines, neither line survives unchanged in the patched body.
Thus:
\[
r_{\text{env}} = \frac{0}{2} = 0.0.
\]
This example shows that recovery is not purely local at the textual level: an early failure may invalidate
later commands that depend on the proof state created by the failed command.

\paragraph{Case 3: Skeleton partial recovery.}
The following skeleton introduces a valid proof step before producing an invalid tactic:
\begin{verbatim}
have h : 1 + 1 = 2 := by
  rfl
apply nonexistent_lemma
exact h
\end{verbatim}
The first two lines form a valid local proof step. The third line fails because the
lemma name cannot be resolved, and the final line becomes unusable after the failed application. The
patched body is:
\begin{verbatim}
have h : 1 + 1 = 2 := by
  rfl
sorry
\end{verbatim}
The original body has four clean lines. The first two lines survive unchanged, while the remaining two
lines are replaced by a placeholder. Therefore:
\[
r_{\text{env}} = \frac{2}{4} = 0.5.
\]
This case illustrates why recovery is useful for skeleton actions: even when the generated proof does not
fully close the goal, valid setup lines and child subgoals can still be preserved and rewarded.

\paragraph{Case 4: Type mismatch.}
A minimal type mismatch example is:
\begin{verbatim}
exact 1
\end{verbatim}
when the expected target is a proposition. Lean rejects the term because a numeral does not inhabit the
expected proposition. The recovered body is:
\begin{verbatim}
sorry
\end{verbatim}
No original clean line survives, giving:
\[
r_{\text{env}} = 0.0.
\]
